\section{Artifact Examples}

This appendix shows concrete artifacts from the live internal-validation run. The examples are generated from the run bundle rather than rewritten by hand.

\subsection{Model Response Examples}

\begin{tabular}{p{0.18\linewidth}p{0.13\linewidth}p{0.61\linewidth}}
\toprule
Model & Dasha cluster & Response excerpt \\
\midrule
gpt-5.2 & original\_\_cluster\_2 & No. \#\#\# 1) The premarital promise is within the Illinois Statute of Frauds and is unenforceable without a writing In Illinois, an oral promise **made in consideration of marriage** (“If you marry me, I will …”) falls within the Statute of Frauds (750 ILCS 80/1, traditional “marriage consideration” provision). A promise to confer a property benefit on a spouse in exchange for marrying—such as a promise to keep the spouse as beneficiary of a death benefit—must be in a signed writing to be enforceable. Here, the promise was explicitly conditioned on marriage and there is **no signed premarital agreement or memorandum**, so the underlying contract claim fails. \#\#\# 2) Marriage (and related rel... \\
gpt-5.2 & original\_\_cluster\_2 & No. \#\#\# 1) The premarital promise is within the Statute of Frauds In Illinois, an oral promise made *in consideration of marriage* (“marry me and I will do X”) falls within the Statute of Frauds and is unenforceable unless there is a writing signed by the promisor (or an equivalent signed memorandum). A promise to name (or keep) a spouse as beneficiary because she marries him is the classic “marriage consideration” contract. So, absent a signed writing, Amanda cannot sue to enforce Clinton’s premarital promise as a contract. \#\#\# 2) Marriage (and related reliance) does not supply part performance to avoid that bar Illinois does not treat the marriage itself as “part performance” that takes... \\
gpt-5.4 & original\_\_cluster\_2 & No. Under Illinois law, Amanda would almost certainly lose, and the later bylaw-compliant change of beneficiary would control. \#\# Why \#\#\# 1. The premarital promise is within the Statute of Frauds A promise made **in consideration of marriage**—for example, “marry me and I’ll give you this property/right/benefit”—must generally be **in writing** to be enforceable, unless it is merely a mutual promise to marry. Clinton’s statement was not just “I’ll marry you.” It was a collateral promise: if Amanda married him, he would make her beneficiary of the fraternal benefit certificate. That is the kind of marriage-based promise Illinois treats as falling within the Statute of Frauds. Because there... \\
gpt-5.4 & original\_\_cluster\_2 & No. Under Illinois law, Amanda would almost certainly lose, and the later bylaw-compliant change of beneficiary would control. \#\# Why \#\#\# 1. The premarital promise is within the Statute of Frauds A promise made **in consideration of marriage**—other than the mutual promise to marry itself—must be in writing to be enforceable. Clinton’s statement was essentially: - *“If you marry me, I will make you beneficiary.”* That is a classic **promise upon consideration of marriage**. Since there is: - no written premarital agreement, and - no signed memorandum, Amanda cannot enforce that promise as a contract. \#\# 2. Marriage itself is not enough part performance Illinois follows the rule that **mar... \\
claude-sonnet-4-6 & original\_\_cluster\_2 & \# Amanda's Claim to the Death Benefit: Legal Analysis \#\# Overview of the Problem This question sits at the intersection of several doctrines: the Statute of Frauds (particularly the marriage provision), promissory estoppel, part performance, the contractual rights of fraternal benefit association members, and the enforceability of beneficiary designations. The short answer is that **Amanda faces serious obstacles and will likely fail**, though the analysis is genuinely complex and some arguments have real force. --- \#\# I. The Threshold Problem: Statute of Frauds \#\#\# A. The Marriage Provision Illinois follows the traditional Statute of Frauds, which requires that promises *made in consider... \\
claude-sonnet-4-6 & original\_\_cluster\_3 & \# Amanda's Claim to the Death Benefit: Legal Analysis \#\# Overview of the Problem This question sits at the intersection of four distinct legal doctrines: (1) the Statute of Frauds as applied to premarital promises, (2) promissory estoppel and part performance, (3) the contractual/quasi-contractual enforceability of beneficiary designations in fraternal benefit certificates, and (4) the internal governance rules of fraternal benefit associations and their legal effect on third-party claims. The short answer is that Amanda faces serious obstacles but has plausible arguments, and the outcome turns substantially on which theory she advances and how Illinois courts weigh competing consideratio... \\
claude-haiku-4-5-20251001 & original\_\_cluster\_4 & \# Analysis of Amanda's Claims to the Death Benefit This is a challenging case involving the intersection of contract law, the Statute of Frauds, promissory estoppel, and the nature of beneficiary designations. Amanda faces significant obstacles, though she has some arguable claims. \#\# I. The Oral Promise and Statute of Frauds \#\#\# The Core Problem Amanda's most direct claim—breach of Clinton's oral promise to name her beneficiary—runs headlong into the **Statute of Frauds**. Under Illinois law (following the Restatement): - **Marriage itself is a consideration** for the promise, so the promise falls within the marriage provision of the Statute of Frauds (contracts made in consideration of... \\
claude-haiku-4-5-20251001 & original\_\_cluster\_4 & \# Analysis of Amanda's Claims This is a nuanced problem involving the intersection of contract law, promissory estoppel, part performance doctrine, and the enforceability of beneficiary designations. Amanda faces substantial obstacles. \#\# I. The Statute of Frauds Problem \#\#\# A. Express Contract Claim Amanda's strongest argument is that Clinton made an enforceable contract: "Marry me and I'll name you beneficiary." However, this faces the **statute of frauds**. **Illinois recognizes a statute of frauds for:** - Contracts in consideration of marriage (Ill. Comp. Stat. 815 ILCS 105/2) - Contracts that cannot be performed within one year **Critical issue:** Is a promise to name someone as a b... \\
google/gemini-3-flash & original\_\_cluster\_2 & Under Illinois law, Amanda likely **cannot** enforce the oral premarital promise to invalidate the later beneficiary change. Her claim would likely fail due to the **Statute of Frauds** and the **Illinois Uniform Premarital Agreement Act (IUPAA)**, and her arguments for equitable relief (part performance or estoppel) are generally unsuccessful in this specific context. Here is the legal breakdown of why Amanda would likely lose: \#\#\# 1. The Writing Requirement (IUPAA and Statute of Frauds) In Illinois, any agreement made in consideration of marriage must be in writing. * **IUPAA (750 ILCS 10/):** Since 1990, Illinois has followed the Uniform Premarital Agreement Act. Section 3 of the Act e... \\
google/gemini-3-flash & original\_\_cluster\_2 & Under Illinois law, Amanda will likely **fail** in her attempt to enforce the oral premarital promise or to invalidate the subsequent beneficiary change. The resolution of this case depends on the intersection of the Illinois Statute of Frauds, the Illinois Uniform Premarital Agreement Act (IUPAA), and the law governing beneficiary designations in fraternal benefit associations. \#\#\# 1. The Statute of Frauds and the IUPAA The primary obstacle for Amanda is the requirement that agreements made in consideration of marriage must be in writing. * **The IUPAA (750 ILCS 10/1 et seq.):** Illinois adopted the Uniform Premarital Agreement Act in 1990. Section 3 of the Act explicitly states: "A prem... \\
google/gemini-3-pro & original\_\_cluster\_2 & Based on Illinois law, **no, Amanda cannot enforce the premarital promise or use estoppel/part performance to invalidate the later beneficiary change.** The death benefit will be awarded to Clinton’s children and sister. Here is the legal analysis of why Amanda’s claim will fail: \#\#\# 1. The Statute of Frauds Under the Illinois Statute of Frauds (740 ILCS 80/1), any agreement made upon consideration of marriage must be in writing and signed by the party to be charged. Clinton’s promise to name Amanda as his beneficiary was made orally and strictly in consideration of her marrying him. Because there is no written premarital agreement or signed memorandum, the contract is unenforceable. \#\#\#... \\
google/gemini-3-pro & original\_\_cluster\_2 & Based on Illinois law, **no, Amanda cannot enforce the premarital promise or use estoppel/part performance to invalidate the later beneficiary change.** The death benefit will be awarded to Clinton’s children and sister. Here is the legal analysis of why Amanda’s claim will fail: \#\#\# 1. The Statute of Frauds Under the Illinois Statute of Frauds (740 ILCS 80/1), any agreement made in consideration of marriage must be in writing and signed by the party to be charged. Because Clinton’s promise (“If you marry me, I will…”) was entirely oral, it is presumptively unenforceable. \#\#\# 2. Marriage is Not "Part Performance" Amanda cannot rely on the act of marriage itself to trigger the "part perfor... \\
anthropic/claude-4.5-sonnet & original\_\_cluster\_3 & This is a complex problem involving contract law, the Statute of Frauds, and the interaction between private agreements and organizational bylaws. Let me work through the key issues: \#\# The Statute of Frauds Problem Clinton's premarital promise faces a significant obstacle: **the Statute of Frauds requires contracts made in consideration of marriage to be in writing.** This is one of the classic categories covered by the Statute of Frauds in Illinois and virtually all jurisdictions. Amanda's theories to overcome this: \#\#\# 1. **Part Performance Exception** - Amanda actually married Clinton in reliance on the promise - Clinton initially performed by designating her as beneficiary - Courts a... \\
anthropic/claude-4.5-sonnet & original\_\_cluster\_3 & This is a complex question involving contract law, the statute of frauds, and the rights of fraternal benefit association members. Let me break down the key issues: \#\# The Premarital Promise and Statute of Frauds **Problem for Amanda**: Clinton's oral premarital promise likely falls within the **statute of frauds** as: - A promise made in consideration of marriage (Illinois requires such promises to be in writing) - Under 740 ILCS 80/1, contracts in consideration of marriage must be in writing **Potential exception**: The **part performance doctrine** might apply: - Amanda fully performed by marrying Clinton - Clinton partially performed by initially naming her beneficiary - However, part... \\
openai/gpt-4.1-mini & original\_\_cluster\_2 & This question involves issues of contract enforceability, the statute of frauds, estoppel/part performance, and the effect of the association’s bylaws on beneficiary designations. --- \#\#\# Facts summary: - Clinton is a member of a fraternal benefit association that pays death benefits to named beneficiaries. - The bylaws allow members to change beneficiaries by surrendering the certificate or, if unable, by affidavit explaining why the certificate cannot be produced. - Clinton orally promises Amanda before marriage that if she marries him, he will surrender his certificate and get a new one naming her beneficiary. - Amanda marries Clinton relying on this promise. - After marriage, Clinton... \\
openai/gpt-4.1-mini & original\_\_cluster\_2 & Let's analyze the situation step-by-step under Illinois law principles related to contract formation, estoppel, and beneficiary designations in fraternal benefit associations. --- \#\#\# Facts Summary: - Clinton has a certificate naming a beneficiary under a fraternal benefit association. - The association's bylaws allow changing beneficiaries either by surrendering the old certificate or by affidavit if surrender is impossible. - Clinton orally promises Amanda before marriage: if she marries him, he will surrender the old certificate and get a new one naming her beneficiary. - Amanda marries Clinton relying on this promise. - After marriage, Clinton obtains a new certificate naming Amanda b... \\
deepseek-ai/deepseek-r1 & original\_\_cluster\_2 & Based on Illinois law regarding fraternal benefit societies, the Statute of Frauds, and promissory estoppel, **Amanda is unlikely to succeed in enforcing the premarital promise or invalidating the later beneficiary change.** The children/sister should prevail. Here's the analysis: 1. **The Statute of Frauds Bars Enforcement of the Contract:** * Clinton's promise ("If you marry me, I will... name you as beneficiary") constitutes a promise made "in consideration of marriage." Under the Illinois Statute of Frauds (740 ILCS 80/1), such promises **must be in writing and signed by the party to be charged** (Clinton) to be enforceable. * There is no written premarital agreement or signed memoran... \\
deepseek-ai/deepseek-r1 & original\_\_cluster\_2 & Based on Illinois law governing fraternal benefit associations and contract principles, **Amanda likely CANNOT enforce the premarital promise or invalidate the later beneficiary change to recover the proceeds.** Here's the breakdown: 1. **The Premarital Promise as a Contract:** * **Statute of Frauds for Premarital Agreements (750 ILCS 10/1 et seq.):** Illinois law requires premarital agreements to be *in writing* and signed by both parties to be enforceable. Clinton's oral promise ("If you marry me, I will... name you as beneficiary") is precisely the type of promise concerning property rights upon marriage that falls within this statute. * **No Signed Writing:** There is no written prema... \\
meta/meta-llama-3-70b-instruct & original\_\_cluster\_1 & A fascinating question! In this scenario, Amanda is relying on the doctrine of promissory estoppel to enforce Clinton's premarital promise. To succeed, Amanda must demonstrate that: 1. Clinton made a clear and unambiguous promise to name her as the beneficiary. 2. Amanda reasonably relied on that promise by marrying Clinton. 3. Her reliance was detrimental, meaning she gave up something of value (in this case, her freedom to marry someone else or not marry at all). 4. Injustice can be avoided only by enforcing the promise. Here, Amanda has a strong case: * Clinton's oral promise was clear and specific. * Amanda reasonably relied on the promise by marrying Clinton, which is a significant c... \\
meta/meta-llama-3-70b-instruct & original\_\_cluster\_1 & What an interesting case! In Illinois, the court would likely apply the doctrine of promissory estoppel to determine whether Amanda can enforce Clinton's premarital promise. Promissory estoppel is a legal doctrine that allows a party to recover on a promise even if it's not supported by consideration (something of value exchanged) if the following elements are met: 1. A promise was made. 2. The promisor (Clinton) should have reasonably expected the promisee (Amanda) to act in reliance on the promise. 3. The promisee did act in reliance on the promise. 4. Injustice can be avoided only by enforcing the promise. In this case: 1. Clinton made an oral promise to name Amanda as the beneficiary... \\
gpt-5.2 & V\_surface\_invariant\_\_cluster\_3 & No. \#\#\# 1) The premarital promise is within the Statute of Frauds and is not enforceable without a writing In Illinois, a promise **made “in consideration of marriage”** (i.e., “marry me and I’ll do X”) falls squarely within the Statute of Frauds and is **unenforceable unless there is a written, signed agreement or memorandum**. Jordan’s statement—“If you marry me, I will surrender my current certificate and obtain a new certificate naming you as beneficiary”—is exactly that kind of promise. \#\#\# 2) Marriage (and related reliance) does not take this kind of agreement out of the Statute of Frauds in Illinois Alex cannot use **part performance** or **equitable estoppel** based on (a) the mar... \\
gpt-5.2 & V\_surface\_invariant\_\_cluster\_3 & No. \#\#\# 1) The premarital promise is within Illinois’s Statute of Frauds and is unenforceable without a writing Illinois requires a signed writing for an **agreement made upon consideration of marriage** (i.e., “marry me and I will do X”), other than mutual promises to marry. Jordan’s statement—“If you marry me, I will … obtain a new certificate naming you as beneficiary”—is exactly that kind of bargain. With **no written premarital agreement or signed memorandum**, Alex cannot enforce the promise as a contract to keep her as beneficiary. \#\#\# 2) Marriage (and the later initial designation) does not take this kind of promise out of the Statute of Frauds Illinois does not treat **marriage i... \\
gpt-5.4 & V\_surface\_invariant\_\_cluster\_3 & No. Under Illinois law, Alex is very unlikely to enforce Jordan’s premarital oral promise or use estoppel/part performance to defeat the later beneficiary change that complied with the association’s bylaws. \#\# Why \#\#\# 1. The oral premarital promise is within the Statute of Frauds A promise made **in consideration of marriage**—“if you marry me, I will do X”—must generally be in writing to be enforceable. Jordan’s promise was exactly that: *marry me and I’ll name you beneficiary.* Because there is **no written premarital agreement or signed memorandum**, Alex cannot enforce that promise as a contract. \#\#\# 2. Marriage itself is not enough part performance Illinois follows the rule that **ma... \\
gpt-5.4 & V\_surface\_invariant\_\_cluster\_3 & No. Under Illinois law, Alex likely cannot enforce Jordan’s premarital oral promise, and the later beneficiary change made in compliance with the association’s bylaws will control. \#\# Why \#\#\# 1. The oral promise is within the Statute of Frauds A promise made **in consideration of marriage**—for example, “if you marry me, I will give you property / make you beneficiary”—must be in writing to be enforceable. Jordan’s statement was exactly that type of promise: **“If you marry me, I will … obtain a new certificate naming you as beneficiary.”** Because there is **no written premarital agreement or signed memorandum**, Alex cannot enforce that promise as a contract. \#\# 2. Marriage itself is no... \\
claude-sonnet-4-6 & V\_surface\_invariant\_\_cluster\_3 & \# Analysis: Alex's Claim Against the Death Benefit Proceeds \#\# Overview of the Legal Problem This question sits at the intersection of several distinct legal doctrines: the Statute of Frauds (particularly the marriage provision), equitable estoppel and part performance, the internal governance rules of fraternal benefit associations, and the competing equities between a surviving spouse and named beneficiaries. The answer is **probably no**—Alex almost certainly cannot prevail—though the analysis is genuinely complex and some equitable arguments deserve serious treatment before being rejected. --- \#\# I. The Threshold Problem: The Statute of Frauds \#\#\# A. The Marriage Provision Illinois fo... \\
claude-sonnet-4-6 & V\_surface\_invariant\_\_cluster\_3 & \# Alex's Claim Against the Death Benefit Proceeds \#\# Executive Summary Alex faces a genuinely difficult case with significant doctrinal obstacles, but has meaningful arguments worth carefully analyzing. The honest assessment is that **Alex likely loses** under orthodox Illinois law, but the claim is not frivolous and a sympathetic court could rule otherwise under equitable theories. Here is the full analysis. --- \#\# I. The Threshold Problem: The Statute of Frauds \#\#\# A. The Promise Falls Within the Marriage Provision Jordan's promise — "marry me and I'll name you beneficiary" — is a promise made **in consideration of marriage**. Under the Illinois Statute of Frauds (740 ILCS 80/1), agreem... \\
claude-haiku-4-5-20251001 & V\_surface\_invariant\_\_cluster\_4 & \# Analysis of Alex's Claim to the Death Benefit This is a challenging case involving the intersection of contract law, the Statute of Frauds, equitable estoppel, and the nature of beneficiary designations. Alex faces significant obstacles but has some arguable claims. \#\# I. The Oral Contract for a Beneficiary Designation \#\#\# A. Statute of Frauds Problem The threshold issue is whether the oral promise is barred by the Statute of Frauds. Illinois courts would likely analyze this under **two potential provisions**: 1. **Marriage provision (Ill. Comp. Stat. 815 ILCS 105/3)**: Promises "in consideration of marriage" must be in writing. Jordan's promise—"If you marry me, I will name you as bene... \\
claude-haiku-4-5-20251001 & V\_surface\_invariant\_\_cluster\_4 & \# Analysis of Alex's Claim to the Death Benefit This is a genuinely difficult question that turns on several intersecting doctrines. The answer is likely **no, Alex cannot recover**, but with meaningful caveats about estoppel theory. \#\# I. The Core Problem: Statute of Frauds \#\#\# A. The Oral Promise and SOF Illinois recognizes that promises made in consideration of marriage are subject to the statute of frauds (Ill. Comp. Stat. 815 ILCS 105/2). The promise—"marry me and I'll name you beneficiary"—is classically within SOF because it's: - Made in consideration of marriage (the marriage itself is the consideration) - Not capable of full performance within one year (the designation is perpetu... \\
google/gemini-3-flash & V\_surface\_invariant\_\_cluster\_3 & Under Illinois law, Alex will likely **not** be able to enforce the oral premarital promise to invalidate the later beneficiary change. The primary obstacles for Alex are the **Statute of Frauds**, the **Illinois Uniform Premarital Agreement Act (IUPAA)**, and the specific way Illinois courts treat marriage as "consideration." Here is the legal breakdown of why Alex’s claim would likely fail: \#\#\# 1. The Statute of Frauds and the IUPAA In Illinois, any agreement made in consideration of marriage must be in writing to be enforceable. * **The Statute of Frauds (740 ILCS 80/1):** This long-standing rule requires that any promise made "upon consideration of marriage" must be in writing and sig... \\
google/gemini-3-flash & V\_surface\_invariant\_\_cluster\_3 & Under Illinois law, Alex will likely **not** be able to enforce the oral premarital promise or invalidate the later beneficiary change. The resolution of this issue depends on the intersection of the Statute of Frauds, the Illinois Uniform Premarital Agreement Act (IUPAA), and the law governing fraternal benefit associations. \#\#\# 1. The Statute of Frauds and the IUPAA In Illinois, agreements made in consideration of marriage are subject to strict writing requirements: * **The Illinois Frauds Act (740 ILCS 80/1):** This statute provides that no action shall be brought to charge any person upon any agreement made upon consideration of marriage, unless the agreement is in writing and signed... \\
google/gemini-3-pro & V\_surface\_invariant\_\_cluster\_2 & Based on Illinois law, **Yes, Alex can recover the proceeds.** Although the initial oral promise was subject to the Statute of Frauds, Jordan’s subsequent act of actually naming Alex as the beneficiary and delivering the certificate to her constituted **full execution** of the contract. This removed the agreement from the Statute of Frauds and gave Alex a **vested equitable interest** in the death benefit that Jordan could not unilaterally defeat. Here is the step-by-step legal analysis of how Illinois courts resolve this exact scenario: \#\#\# 1. The Statute of Frauds and "Part Performance" Under the Illinois Statute of Frauds (740 ILCS 80/1), any agreement made in consideration of marriage... \\
google/gemini-3-pro & V\_surface\_invariant\_\_cluster\_3 & Based on Illinois law, **Alex cannot enforce the premarital promise or use estoppel/part performance to invalidate the later beneficiary change.** The death benefit will legally go to Jordan’s children and sister. Here is the step-by-step legal analysis of why Alex’s claims will fail: \#\#\# 1. The Nature of the Beneficiary’s Interest Under Illinois law governing fraternal benefit associations and life insurance, a named beneficiary generally has a "mere expectancy" rather than a vested right to the proceeds. The member retains the absolute right to change the beneficiary at any time, provided they comply with the association's constitution and bylaws. The only exception to this rule is if t... \\
anthropic/claude-4.5-sonnet & V\_surface\_invariant\_\_cluster\_3 & This is an excellent multi-layered question involving contract law, the statute of frauds, and the intersection of promissory estoppel with association governance rules. \#\# Alex's Likely Outcome: **Cannot Recover** Despite some equitable arguments, Alex will likely lose because: \#\#\# 1. **The Premarital Promise Issue is Moot** Alex doesn't need to rely on the oral premarital promise because **Jordan actually performed it**—he did change the beneficiary designation to Alex after marriage. The promise was fulfilled, so statute of frauds issues (which would apply to premarital promises in many jurisdictions) become irrelevant. \#\#\# 2. **Alex Received a Revocable Designation** The critical prob... \\
anthropic/claude-4.5-sonnet & V\_surface\_invariant\_\_cluster\_4 & This is a complex question involving the intersection of contract law, the statute of frauds, promissory estoppel, and the internal rules of a fraternal benefit association. Here's my analysis: \#\# The Core Legal Issues \#\#\# 1. **Statute of Frauds Problem** Jordan's premarital promise likely falls under the statute of frauds as: - A **promise made in consideration of marriage** (one of the classic categories) - Illinois, like most states, requires such promises to be in writing **Alex's problem**: The oral promise alone is likely unenforceable under the statute of frauds. \#\#\# 2. **Part Performance/Promissory Estoppel Arguments** Alex has stronger arguments here: **Part Performance Exception... \\
openai/gpt-4.1-mini & V\_surface\_invariant\_\_cluster\_3 & This is a question involving contract enforceability, estoppel, and the effect of the fraternal benefit association’s bylaws on beneficiary designations. --- \#\#\# Facts Summary: - Jordan was a member of a fraternal benefit association that pays death benefits to designated beneficiaries. - The association’s bylaws allow changing beneficiaries by surrendering the old certificate or, if the certificate cannot be surrendered, by an affidavit explaining why and naming a new beneficiary. - Jordan orally promised Alex before marriage that he would surrender the old certificate and name Alex as beneficiary if she married him. - Alex married Jordan relying on that promise. - After marriage, Jordan... \\
openai/gpt-4.1-mini & V\_surface\_invariant\_\_cluster\_3 & This question involves Illinois contract and insurance law principles, focusing on: - The enforceability of an oral premarital promise concerning beneficiary designation, - Whether estoppel or part performance can validate that promise despite the lack of a writing, - The effect of the association’s bylaws on beneficiary changes, - And the competing claims to the death benefit after the last beneficiary change. --- \#\#\# Key Facts and Issues: - Jordan was a long-time member of a fraternal benefit association. - The association’s constitution/bylaws allow a member to change beneficiaries via surrender of certificate or affidavit if surrender is impossible. - Jordan promised Alex orally befor... \\
deepseek-ai/deepseek-r1 & V\_surface\_invariant\_\_cluster\_3 & Based on Illinois law, **Alex cannot enforce Jordan's premarital promise or use estoppel/part performance to invalidate the later beneficiary change.** The children/sister are likely entitled to the death benefit. Here's the analysis: 1. **Statute of Frauds Bars Enforcement of the Premarital Promise:** * Jordan's oral promise ("If you marry me, I will... name you as beneficiary") is a promise to make a testamentary disposition (a disposition taking effect at death) or a conditional gift in contemplation of marriage. * Illinois' Statute of Frauds (740 ILCS 80/1, derived from the English Statute of Frauds) **requires promises to make testamentary dispositions or gifts in consideration of ma... \\
deepseek-ai/deepseek-r1 & V\_surface\_invariant\_\_cluster\_3 & Based on Illinois law, **Alex cannot enforce the premarital promise or use estoppel/part performance to invalidate the later beneficiary change.** The children/sister are entitled to the death benefit. Here's the analysis: 1. **Unenforceability of the Premarital Promise:** * Jordan's oral promise ("marry me and I'll name you beneficiary") is a promise to make a future disposition of property (the death benefit) in consideration of marriage. * Under Illinois law (the **Statute of Frauds**), such agreements must be **in writing and signed by the party to be charged** (Jordan) to be enforceable. *See 740 ILCS 80/1; In re Estate of Sewart*, 274 Ill. App. 3d 298 (1995). * Since there's no writ... \\
meta/meta-llama-3-70b-instruct & V\_surface\_invariant\_\_cluster\_1 & A fascinating question! In this scenario, Alex is arguing that Jordan's premarital promise to name her as beneficiary should be enforced, despite the later change in beneficiaries. To succeed, Alex would need to rely on the doctrine of promissory estoppel or part performance. Promissory estoppel is a legal doctrine that allows a party to recover benefits conferred in reliance on a promise, even if there is no formal contract. To establish promissory estoppel, Alex would need to show that: 1. Jordan made a clear and unambiguous promise to name her as beneficiary. 2. Alex reasonably relied on that promise by marrying Jordan. 3. Alex suffered a detriment (i.e., she married Jordan) in relianc... \\
meta/meta-llama-3-70b-instruct & V\_surface\_invariant\_\_cluster\_5 & What an interesting scenario! In this case, Alex is arguing that Jordan's premarital promise to name her as beneficiary should be enforced, despite the lack of a written agreement. Alex is relying on the doctrine of promissory estoppel, which allows a party to recover damages or enforce a promise when they have reasonably relied on that promise to their detriment. To succeed, Alex must show that: 1. Jordan made a clear and unambiguous promise to name her as beneficiary. 2. Alex reasonably relied on that promise by marrying Jordan. 3. Alex's reliance was to her detriment (i.e., she would not have married Jordan without the promise). 4. Jordan's later attempt to change beneficiaries is unco... \\
gpt-5.2 & V1\_\_cluster\_6 & Yes. The signed premarital note is enough (as a Statute of Frauds matter) to make the *marriage‑consideration contract* enforceable, and that gives Amanda a basis to attack the later beneficiary change—though her success still turns on beneficiary/insurance rules and the association’s change‑of‑beneficiary provisions. Under Illinois Statute of Frauds principles (as commonly stated): - **A promise “in consideration of marriage” must be in a signed writing to be enforceable.** Clinton’s dated, signed note—“In consideration of Amanda marrying me, I agree to name her beneficiary…”—fits that category. It identifies the parties, states the marriage consideration, and states the essential promis... \\
gpt-5.2 & V1\_\_cluster\_3 & Yes. Under Illinois Statute of Frauds principles, the signed premarital writing is what makes Amanda’s “marriage-consideration” agreement enforceable, and that enforceability is what gives her a basis to attack the later beneficiary change. - **Agreement “in consideration of marriage” must be in writing and signed.** A promise made in consideration of marriage (other than mutual promises to marry) falls within the Statute of Frauds. Clinton’s signed, dated note saying, in substance, “in consideration of Amanda marrying me, I agree to name her beneficiary” is the kind of memorandum that satisfies that requirement. So **Amanda can enforce the agreement**. - **If there were no signed writing... \\
gpt-5.4 & V1\_\_cluster\_1 & Yes. The signed premarital writing is the key fact that can take Amanda’s marriage-consideration agreement out of the Statute of Frauds problem. \#\#\# Why An agreement made **in consideration of marriage** generally falls within the **Statute of Frauds** and must be supported by a **writing signed by the party to be charged**. Here, before the marriage, Clinton signed a note stating: > “In consideration of Amanda marrying me, I agree to name her beneficiary on my association certificate.” That is strong evidence of a signed written memorandum of the agreement. So, unlike a purely oral antenuptial promise, this promise is likely **enforceable** under Statute of Frauds principles. \#\#\# Effect... \\
gpt-5.4 & V1\_\_cluster\_1 & Yes. The signed premarital writing is highly important, because it likely satisfies the Statute of Frauds requirement for an agreement made in consideration of marriage. \#\# Why A promise made **in consideration of marriage** generally falls within the Statute of Frauds and must be **in writing and signed by the party to be charged**. Here, before the marriage, Clinton signed and dated a writing stating: > “In consideration of Amanda marrying me, I agree to name her beneficiary on my association certificate.” That is strong evidence of an enforceable marriage-consideration contract, because: - it is **written**, - it is **signed by Clinton**, - it identifies the **consideration** (Amanda m... \\
claude-sonnet-4-6 & V1\_\_cluster\_6 & \# Analysis of the Premarital Writing Under Statute of Frauds Principles \#\# The Core Statute of Frauds Issue The critical question is whether Clinton's promise to name Amanda as beneficiary — made **in consideration of marriage** — falls within the Statute of Frauds, which in most jurisdictions (including Illinois) requires that agreements made **in consideration of marriage** be in writing to be enforceable. --- \#\# The Writing Clinton Signed Clinton's note stated: > *"In consideration of Amanda marrying me, I agree to name her beneficiary on my association certificate."* This document must be evaluated against Statute of Frauds requirements for sufficiency of a writing. A memorandum satis... \\
claude-sonnet-4-6 & V1\_\_cluster\_6 & \# Analysis of the Premarital Writing Under Statute of Frauds Principles \#\# The Core Issue The question is whether Clinton's signed premarital note — promising to name Amanda as beneficiary **in consideration of her marrying him** — satisfies Statute of Frauds requirements and thereby strengthens Amanda's ability to enforce the agreement and challenge the subsequent beneficiary change. --- \#\# The Relevant Statute of Frauds Concern Under traditional Statute of Frauds principles, a **promise made in consideration of marriage** (other than a mutual promise to marry) must generally be **in writing and signed** by the party to be charged to be enforceable. This is the so-called **"marriage cons... \\
claude-haiku-4-5-20251001 & V1\_\_cluster\_1 & \# Analysis Under Illinois Statute of Frauds This question involves the intersection of the Statute of Frauds, marriage consideration, and beneficiary designation law. Here's the analysis: \#\# The Statute of Frauds Issue Under Illinois law, an agreement that is made "in consideration of marriage" must be in writing to be enforceable—this is a classic Statute of Frauds category. **Clinton's signed, dated note satisfies this requirement.** The writing: - Is signed by Clinton (the party to be charged) - Is dated - Clearly expresses the agreement - Explicitly references marriage consideration ("In consideration of Amanda marrying me") \#\# The Marriage Consideration Doctrine Illinois recognizes t... \\
claude-haiku-4-5-20251001 & V1\_\_cluster\_4 & \# Analysis Under Illinois Statute of Frauds This question requires examining whether the premarital writing satisfies the Statute of Frauds and whether it affects Amanda's enforceability claims. \#\# The Statute of Frauds Issue Under Illinois law, agreements made in consideration of marriage must satisfy the Statute of Frauds and be in writing. The signed premarital note stating Clinton's agreement to name Amanda as beneficiary appears to satisfy this requirement because: 1. **Writing exists**: Clinton signed and dated a written note 2. **Clear terms**: It explicitly states the agreement and identifies the consideration (Amanda's marriage) 3. **Signature**: Clinton's signature is present \#\#... \\
google/gemini-3-flash & V1\_\_cluster\_6 & Under Illinois law and general principles of equity regarding fraternal benefit associations, the signed premarital writing **significantly affects** the outcome by satisfying the Statute of Frauds and establishing that Amanda has an **equitable vested interest** in the death benefit. The analysis breaks down as follows: \#\#\# 1. Satisfaction of the Statute of Frauds In Illinois, the Statute of Frauds (740 ILCS 80/1) provides that any agreement made upon consideration of marriage must be in writing and signed by the party to be charged. * **The Writing:** Clinton signed and dated a note stating, “In consideration of Amanda marrying me, I agree to name her beneficiary on my association certi... \\
google/gemini-3-flash & V1\_\_cluster\_5 & Under Illinois law and general principles of insurance equity, the signed premarital writing **does affect the outcome by satisfying the Statute of Frauds, thereby allowing Amanda to enforce the agreement and successfully challenge the later beneficiary change.** Here is the legal breakdown of why the writing is decisive in this scenario: \#\#\# 1. Satisfaction of the Statute of Frauds In Illinois, the Statute of Frauds (740 ILCS 80/1) requires that any agreement made "upon consideration of marriage" must be in writing and signed by the party to be charged. * **The Writing:** Clinton’s note specifically states the consideration ("In consideration of Amanda marrying me") and the promise ("I a... \\
google/gemini-3-pro & V1\_\_cluster\_6 & **Yes, the signed premarital writing fundamentally affects the outcome of this dispute.** Because of the signed writing, Amanda will be able to enforce the agreement, successfully challenge the later beneficiary change, and claim the death benefit. Here is how the Statute of Frauds and Illinois equitable principles apply to these facts: \#\#\# 1. Satisfaction of the Statute of Frauds Under the Illinois Statute of Frauds (740 ILCS 80/1), any agreement made "in consideration of marriage" must be in writing and signed by the party to be charged to be legally enforceable. * **The Problem without a Writing:** If Clinton had only made an oral promise to name Amanda as his beneficiary in exchange f... \\
google/gemini-3-pro & V1\_\_cluster\_5 & **Yes.** The signed premarital writing fundamentally affects Amanda’s rights. Because the writing satisfies the Statute of Frauds, Amanda can successfully enforce the agreement, establish a vested equitable interest in the certificate, and defeat the claims of the later-named beneficiaries. Here is a breakdown of how the legal principles apply to this specific scenario under Illinois law: \#\#\# 1. Satisfaction of the Statute of Frauds Under the Illinois Statute of Frauds, any agreement made in consideration of marriage must be in writing and signed by the party to be charged to be legally enforceable. * **Without the writing:** If Clinton had only made an oral promise to name Amanda as bene... \\
anthropic/claude-4.5-sonnet & V1\_\_cluster\_1 & \# Analysis of the Premarital Writing Under Statute of Frauds \#\# Does the Writing Satisfy the Statute of Frauds? **Yes, the signed premarital writing likely satisfies the Statute of Frauds** for the marriage-consideration promise, which affects Amanda's ability to enforce the agreement. \#\# Key Elements Present \#\#\# 1. **Writing Requirement - SATISFIED** - Clinton signed and dated a written note - The writing memorializes the essential terms of the agreement \#\#\# 2. **Marriage Consideration Exception - APPLICABLE** - The promise is "in consideration of marriage" - This falls within the traditional Statute of Frauds categories requiring a writing - Illinois follows this requirement \#\#\# 3. **Es... \\
anthropic/claude-4.5-sonnet & V1\_\_cluster\_1 & \# Analysis of the Premarital Writing Under Statute of Frauds \#\# The Core Issue The question asks whether Clinton's signed premarital note satisfies the Statute of Frauds for a marriage-consideration contract, potentially giving Amanda grounds to challenge the later beneficiary change. \#\# Statute of Frauds Requirements Under the Statute of Frauds, **contracts in consideration of marriage** must be in writing. The key requirements are: 1. A writing signed by the party to be charged (Clinton) 2. Containing essential terms 3. Showing consideration of marriage \#\# Analysis of Clinton's Note **What the writing provides:** - ✓ Signed and dated by Clinton (the promisor) - ✓ Explicitly states marri... \\
openai/gpt-4.1-mini & V1\_\_cluster\_2 & Let's analyze the situation step-by-step, focusing on Illinois law regarding beneficiary designations, contract enforceability, and the Statute of Frauds. --- \#\#\# Facts Summary: - Clinton is a member of a fraternal benefit association. - The association's bylaws allow Clinton to change beneficiaries, including by affidavit if the original certificate cannot be surrendered. - Before marrying Amanda, Clinton signs a short written note: *“In consideration of Amanda marrying me, I agree to name her beneficiary on my association certificate.”* - Amanda marries Clinton. - After marriage, Clinton obtains a certificate naming Amanda as beneficiary, which Amanda keeps. - One year later, Clinton cl... \\
openai/gpt-4.1-mini & V1\_\_cluster\_2 & Let's analyze the facts and the legal issues step-by-step: --- \#\#\# FACTS: - Clinton is a member of a fraternal benefit association. - The association's bylaws allow Clinton to change beneficiaries by: 1. Direct surrender of the certificate and re-issuance, or 2. By affidavit if he cannot produce the certificate. - Before marrying Amanda, Clinton signs and dates a short written note stating: *“In consideration of Amanda marrying me, I agree to name her beneficiary on my association certificate.”* - Amanda marries Clinton. - After marriage, the association issues a certificate naming Amanda as beneficiary, and Amanda keeps it. - A year later, Clinton claims he cannot produce the certificate... \\
deepseek-ai/deepseek-r1 & V1\_\_cluster\_6 & The signed premarital writing **does affect** Amanda's ability to enforce the marriage-consideration agreement and challenge the later beneficiary change under Statute of Frauds principles. Here's why: 1. **Satisfaction of Statute of Frauds:** The premarital writing signed by Clinton ("In consideration of Amanda marrying me, I agree to name her beneficiary on my association certificate") **satisfies the Statute of Frauds requirement** for an agreement made in consideration of marriage. * It is **in writing**. * It clearly **identifies the subject matter** (naming Amanda beneficiary of the association certificate). * It **states the consideration** (Amanda marrying Clinton). * It is **sign... \\
deepseek-ai/deepseek-r1 & V1\_\_cluster\_1 & Based on the Statute of Frauds principles described: 1. **The Premarital Writing Satisfies the Statute of Frauds:** The signed, dated note stating "In consideration of Amanda marrying me, I agree to name her beneficiary on my association certificate" constitutes a sufficient writing to satisfy the Statute of Frauds for a contract made in consideration of marriage. It is in writing, signed by the party to be charged (Clinton), and identifies the consideration (marriage) and the subject matter (naming Amanda beneficiary). 2. **The Agreement is Enforceable:** Therefore, the agreement between Clinton and Amanda is enforceable. Amanda performed her obligation by marrying Clinton. Clinton parti... \\
meta/meta-llama-3-70b-instruct & V1\_\_cluster\_6 & A fascinating question! Under the Statute of Frauds principles, a contract or agreement that falls within the statute's scope must be in writing and signed by the party to be charged (in this case, Clinton). The premarital writing signed by Clinton, which states "In consideration of Amanda marrying me, I agree to name her beneficiary on my association certificate," can be considered a written agreement that satisfies the Statute of Frauds. This writing constitutes a promise to perform a specific act (naming Amanda as beneficiary) in exchange for Amanda's performance of marrying Clinton. The writing is signed and dated, which meets the formal requirements of the Statute of Frauds. As a res... \\
meta/meta-llama-3-70b-instruct & V1\_\_cluster\_6 & A classic contract law problem! Under the Statute of Frauds, a contract made in consideration of marriage must be in writing and signed by the party to be charged (in this case, Clinton) to be enforceable. The premarital writing signed by Clinton, "In consideration of Amanda marrying me, I agree to name her beneficiary on my association certificate," satisfies the Statute of Frauds requirements. This writing is a contractual agreement made in consideration of marriage, and it is in writing, signed, and dated. As a result, Amanda can enforce this agreement and challenge the later beneficiary change. The fact that Clinton later signed an affidavit and changed the beneficiaries does not nega... \\
\bottomrule
\end{tabular}

\subsection{Dasha Cluster Examples}

\begin{tabular}{p{0.12\linewidth}p{0.22\linewidth}p{0.22\linewidth}p{0.34\linewidth}}
\toprule
Cluster & Member models & Outcome & Reasoning signature excerpt \\
\midrule
original\_\_cluster\_1 & meta/meta-llama-3-70b-instruct (n=2) & Enforce the premarital promise via promissory estoppel; invalidate the later beneficiary change; pay the death benefit to Amanda. & Treat the case as promissory estoppel: Clinton made a clear promise to name Amanda beneficiary; Amanda reasonably relied by marrying; marriage is a substantial detrimental change; injustice requires enforcement; bylaws allowing beneficiary changes do not defeat equitable enforcement. \\
original\_\_cluster\_2 & claude-sonnet-4-6, deepseek-ai/deepseek-r1 (n=2), google/gemini-3-flash (n=2), google/gemini-3-pro (n=2), gpt-5.2 (n=2), gpt-5.4 (n=2), openai/gpt-4.1-mini (n=2) & No enforceable contract right: the oral marriage-consideration promise is barred by the Statute of Frauds; marriage/initial designation does not remove the bar; a later beneficiary change made pursuant to bylaws contr... & Because the promise was made in consideration of marriage, it falls within the Statute of Frauds; there is no signed writing/memorandum; marriage (the contemplated consideration) and the initial post-marriage designation do not constitute part performance to avoid the statute; therefore the promise is a nullity and does not bar the member from later changing beneficiaries under the association bylaws, so the last valid designation governs. \\
original\_\_cluster\_3 & anthropic/claude-4.5-sonnet (n=2), claude-sonnet-4-6 & Amanda likely cannot enforce the premarital oral promise; absent a sufficient writing, Clinton could change beneficiaries pursuant to the association bylaws and the later-named beneficiaries take. & The promise was made in consideration of marriage and thus falls within the Statute of Frauds; there is no signed writing/memorandum, so the premarital promise is unenforceable. Subsequent marriage (and even post-marriage initial designation) does not take it out of the statute; the designation is treated as voluntary and, under the fraternal association contract/bylaws reserving a change right, Clinton remained free to change beneficiaries. \\
original\_\_cluster\_4 & claude-haiku-4-5-20251001 (n=2) & The oral premarital bargain is barred by the Statute of Frauds absent a signed writing; the response nonetheless treats equitable enforcement as possible but uncertain (especially via promissory estoppel). & Because the promise was made in consideration of marriage, it triggers the Statute of Frauds; there is no signed memorandum; therefore the breach-of-contract claim fails. The response then pivots to a potential (but uncertain) promissory-estoppel-based enforcement theory to avoid injustice, while acknowledging bylaws make beneficiary designations revocable and may render reliance unreasonable. \\
V\_surface\_invariant\_\_cluster\_1 & meta/meta-llama-3-70b-instruct & Enforce the premarital promise via promissory estoppel/part performance and invalidate or estop recognition of the later beneficiary change (award benefit to the spouse). & Treat the case as one for equitable enforcement: the insured made a clear premarital promise; the spouse reasonably relied by marrying; marriage and initial designation constitute (partial) performance; injustice warrants estopping denial of the promise notwithstanding bylaws. \\
V\_surface\_invariant\_\_cluster\_2 & google/gemini-3-pro & Alex recovers the proceeds; later-named beneficiaries are subordinated. & The promise is within the marriage-consideration Statute of Frauds, but the insured’s subsequent procurement/issuance and delivery of the beneficiary certificate fully executed the bargain, taking it outside the Statute of Frauds; execution plus marriage consideration gives the spouse a vested equitable interest that defeats a later bylaw-compliant change to volunteer beneficiaries. \\
V\_surface\_invariant\_\_cluster\_3 & anthropic/claude-4.5-sonnet, claude-sonnet-4-6 (n=2), deepseek-ai/deepseek-r1 (n=2), google/gemini-3-flash (n=2), google/gemini-3-pro, gpt-5.2 (n=2), gpt-5.4 (n=2), openai/gpt-4... & The premarital oral promise is unenforceable absent a signed writing; marriage/reliance does not remove the bar; the later bylaw-compliant beneficiary change governs and the later-named beneficiaries take. & Treat the premarital bargain as an agreement in consideration of marriage (not merely a voluntary designation) => Statute of Frauds applies and requires a signed writing => no writing so the promise/contract-not-to-change is a nullity => subsequent marriage and initial designation do not constitute part performance/estoppel to enforce it => absent an enforceable restriction, beneficiary interests remain revocable and the member's later bylaw-compliant change controls the payout. \\
V\_surface\_invariant\_\_cluster\_4 & anthropic/claude-4.5-sonnet, claude-haiku-4-5-20251001 (n=2) & Oral promise is likely unenforceable; part performance is unlikely to remove it from the statute; later beneficiary change under bylaws likely stands (with only a tentative possibility of equitable-estoppel/damages ag... & Classify the promise as a marriage-consideration agreement (not merely a voluntary beneficiary designation), triggering a writing requirement; because there is no signed writing, the promise is not directly enforceable and does not invalidate a later bylaw-compliant beneficiary change. \\
V\_surface\_invariant\_\_cluster\_5 & meta/meta-llama-3-70b-instruct & Uncertain / conditional: court could enforce via promissory estoppel if reliance and injustice are shown; otherwise it will likely honor the bylaws and allow the later beneficiary change. & Treats the claim as promissory estoppel: identifies elements (clear promise, reasonable reliance by marrying, detriment, injustice/unconscionability) and concludes enforcement is possible if those elements are satisfied, potentially invalidating the later beneficiary change. \\
V1\_\_cluster\_1 & anthropic/claude-4.5-sonnet (n=2), claude-haiku-4-5-20251001, deepseek-ai/deepseek-r1, gpt-5.4 (n=2) & Enforceable: the signed premarital writing likely satisfies the Statute of Frauds, strengthening the spouse’s claim and supporting a challenge to the later beneficiary change as a breach of the enforceable agreement. & G1 satisfied because the promise is expressly in consideration of marriage; G2 satisfied because the premarital note is a signed written memorandum by Clinton stating the bargain (marriage in exchange for naming Amanda beneficiary). With the Statute of Frauds satisfied, the promise is enforceable and can be used to contest the later beneficiary change as breach. \\
V1\_\_cluster\_2 & openai/gpt-4.1-mini (n=2) & Enforceable premarital agreement; spouse can enforce and defeat the later beneficiary change; spouse should receive the death benefit. & Treats the premarital note as a sufficient signed memorandum satisfying the Statute of Frauds for a marriage-consideration contract to name the spouse beneficiary; characterizes marriage as consideration; concludes the contract restricts/disallows later change away from spouse absent consent, so later beneficiaries lose. \\
V1\_\_cluster\_3 & gpt-5.2 & Yes—because there is a signed premarital writing/memorandum, the marriage-consideration agreement is enforceable and can support attacking the later beneficiary change (potentially via equitable relief). & Classify the promise as an agreement in consideration of marriage (not merely mutual promises) (G1) → find a signed premarital note as sufficient memorandum satisfying the writing requirement (G2) → conclude the agreement is enforceable and can be used to argue the insured was contractually bound not to change the beneficiary, supporting a challenge to the later bylaw-compliant change (often via constructive trust). \\
V1\_\_cluster\_4 & claude-haiku-4-5-20251001 & Treats the signed premarital note as satisfying the Statute of Frauds, making the agreement enforceable and strengthening Amanda’s ability to challenge the later beneficiary change as a breach. & Because the promise was made in consideration of marriage, the Statute of Frauds applies; the signed, dated note memorializes the essential terms and consideration and is signed by Clinton, so the writing requirement is met and the agreement is enforceable, supporting a claim that later beneficiary changes breach the obligation. \\
V1\_\_cluster\_5 & google/gemini-3-flash, google/gemini-3-pro & The signed premarital writing satisfies the Statute of Frauds, making the agreement enforceable; Amanda can enforce the promise and defeat the later beneficiary change. & Because the promise was in consideration of marriage, the Statute of Frauds applies; the signed note is a sufficient memorandum identifying consideration and promise, so the contract is enforceable. Upon marriage, the spouse performs and acquires an equitable vested interest; later beneficiaries are volunteers and cannot take against the contractual/equitable claim, so the bylaw-based change is ineffective as against Amanda. \\
V1\_\_cluster\_6 & claude-sonnet-4-6 (n=2), deepseek-ai/deepseek-r1, google/gemini-3-flash, google/gemini-3-pro, gpt-5.2, meta/meta-llama-3-70b-instruct (n=2) & Enforceable as a Statute of Frauds matter because the signed premarital writing is treated as a sufficient memorandum; spouse has a basis to contest the later beneficiary change (subject to beneficiary/bylaw rules). & Treat the promise as a marriage-consideration contract within the SOF; find a signed note that states essential terms and thus satisfies the writing requirement; conclude the contract is enforceable and can support a breach/constructive-trust style challenge to the later beneficiary change even if the association followed its procedures. \\
\bottomrule
\end{tabular}
